\chapter*{Preface}
\addcontentsline{toc}{chapter}{Preface}

There is a moment, somewhere past your first algorithms course, when graphs stop
being a topic and start being a language. You stop asking ``how do I run BFS''
and start asking ``what is the right graph, and what does its structure already
know about my problem.'' A scheduling constraint becomes a directed edge; a
bottleneck becomes a cut; a notion of similarity becomes a weight; a hard
combinatorial question dissolves into an eigenvalue. This book is written for the
reader who wants to make that crossing deliberately, and to come out the other
side able to prove that the tools work and to say exactly how fast they run.

The title is a promise about altitude. We start on the ground, with breadth-first
search on a tiny graph you can draw on a napkin, and we climb -- without ever
asking you to take a leap of faith -- all the way to Laplacian solvers that
compute electrical flows in nearly linear time, machinery that sits at the
research frontier and powers the fastest known algorithms for problems as old as
maximum flow. The climb is long, but every step is a step, never a jump.

\chapterline

\section*{Who this book is for}

This is the \emph{advanced companion} to a friendly introductory algorithms book.
We assume you have met the classics before, even if the memory is hazy: you have
seen BFS and DFS, you know what a stack and a queue are, you can read a recursion,
you have heard the words ``Dijkstra'' and ``minimum spanning tree'' and have a
rough sense of what they do. We assume you are comfortable with big-$\Oh$ notation
as an idea and will not flinch when we write $\Th(m + n\log n)$.

What we do \emph{not} assume is that you remember the proofs, or that you have
ever seen anything past the introductory layer. We re-derive the foundations at a
faster pace and a higher rigor than an intro course, precisely so that the
advanced material -- blossoms, link-cut trees, treewidth, Cheeger's inequality,
electrical flows -- rests on bedrock and not on half-remembered intuition. If you
can comfortably follow a proof by induction and you are willing to draw the
occasional picture, you are ready.

\begin{intuitionbox}
A useful self-test. Take a connected undirected graph and remove one edge. When
does the graph fall into two pieces? If the phrase ``that edge was a bridge''
came to mind, you have the right background. If it did not, you will after the
chapter on bridges and biconnectivity -- but you should expect to work for it.
This book rewards the reader who pauses, draws, and re-derives.
\end{intuitionbox}

\section*{What ``from the ground up'' actually means here}

Many advanced books open in the middle. They state a theorem, prove it cleanly,
and move on, leaving you to reconstruct the path that a researcher actually walked
to arrive there. We do the opposite. For every major algorithm the chapter follows
the same arc, and the arc is the pedagogy:

\begin{enumerate}
  \item \textbf{A picture and an intuition.} Before a single line of pseudocode we
    draw a small graph and ask what \emph{should} happen, building the mental model
    that the algorithm will later make precise. Intuition first is not a courtesy;
    it is how you will later \emph{remember} the algorithm and how you will debug
    it when your implementation misbehaves at 2~a.m.
  \item \textbf{The algorithm, in clean pseudocode.} Every major algorithm is
    stated in an \texttt{algorithm2e} block precise enough to implement directly,
    with the data structures named explicitly.
  \item \textbf{A correctness argument.} We prove the thing works -- via a loop
    invariant, an exchange argument, an augmenting-path argument, a duality, or a
    potential function. The proof technique is itself a transferable skill, and we
    name it each time so the catalogue grows in your head.
  \item \textbf{A tight complexity analysis.} Not ``it's polynomial'' but the
    exact bound, with the bottleneck identified and, where it matters, a matching
    lower bound so you know the analysis is not loose.
  \item \textbf{A trace by hand.} We run the algorithm on a concrete small graph,
    showing the data-structure state at each step. This is the single most
    important habit the book tries to instill.
\end{enumerate}

\begin{whybox}
Why insist on intuition \emph{and} full rigor, rather than picking one? Because
they fail in opposite directions, and the pair is robust where either alone is
brittle. Pure intuition produces algorithms that are ``intuitively plausible'' and
quietly wrong on the adversarial input -- the negative-weight cycle, the
zero-capacity edge, the disconnected component. Pure formalism produces proofs you
can verify line by line but cannot reconstruct, modify, or apply to the next
problem. Hold both and you get an algorithm you can re-derive from scratch on a
whiteboard \emph{and} defend against an adversary. That combination is what
separates someone who has read about algorithms from someone who can build them.
\end{whybox}

\section*{The colored-box language}

Throughout the book a small set of recurring boxes carries the load that a good
lecturer carries with tone of voice -- this part is the heart, this part is a
landmine, this is a war story. Learn to read the colors and you can skim or dive
at will.

\begin{focusbox}
A \textbf{Focus} box frames the one question a section is really about. When you
are lost, scroll up to the nearest Focus box; it tells you what the surrounding
machinery is \emph{for}.
\end{focusbox}

\begin{intuitionbox}
An \textbf{Intuition} box is the picture-in-words: the mental model, the analogy,
the ``what should happen'' that precedes the formal statement. These are safe to
read first and out of order.
\end{intuitionbox}

\begin{whybox}
A \textbf{Why it works} box is the load-bearing idea of a correctness proof,
stated before the proof so you read the proof knowing where it is going.
\end{whybox}

\begin{trapbox}
A \textbf{Common trap} box is the subtle bug -- the off-by-one, the case the proof
quietly assumes away, the input that breaks the naive version. We hunt these on
purpose. Most of what you will get wrong in practice is collected here in advance.
\end{trapbox}

A \textbf{History} box places an idea in time: Dijkstra at a cafe table in 1956,
Edmonds inventing ``good characterization'' and blossoms, Tarjan turning
depth-first search into a versatile search method, Karger randomizing the minimum cut,
Spielman and Teng solving Laplacian systems in near-linear time. Algorithms are made by
people solving specific problems, and knowing the context clarifies why the algorithm was built.

\begin{complexitybox}
A \textbf{Cost summary} box collects the running time and space of an algorithm in
one place, often with the bottleneck circled and variants compared. When you are
choosing an algorithm for real work, these are the boxes you want.
\end{complexitybox}

\begin{labbox}
A \textbf{Trace it yourself} box hands you a specific small graph and asks you to
run the algorithm by hand before reading our trace. Do these. They are where the
understanding actually happens; reading our trace afterward only confirms it.
\end{labbox}

\begin{applicationbox}
An \textbf{Application} box shows where the algorithm is used in practice --
register allocation, image segmentation, recommendation, routing, circuit
analysis. Abstractions stick better when anchored to a concrete problem.
\end{applicationbox}

\begin{connectionbox}
A \textbf{Connection} box is the most important box for this companion, because it
is the bridge between the three books. It links forward and backward across
chapters here, and directly out to the companion \emph{intro algorithms}
book and the companion \emph{linear algebra} book. When a graph Laplacian turns
out to be a positive semidefinite matrix whose spectrum you already studied, a
Connection box says so explicitly.
\end{connectionbox}

There are a few more -- \emph{Formal} for the precise statement, \emph{Walking
question} for something to chew on, \emph{Pause} to recover and consolidate -- but
the seven above are the spine.

\section*{How the three books fit together}

This volume is one of a trio, and the trio is engineered so that the seams line
up.

\begin{connectionbox}
\textbf{The intro algorithms book} is the prequel. It teaches you to read
pseudocode, to reason about big-$\Oh$, to wield recursion and a handful of core
data structures, and it introduces BFS, DFS, Dijkstra, and union-find at a gentle
pace. Where that book \emph{introduces} a graph algorithm, this book \emph{proves
it tight and then goes far past it}. If a derivation here feels too fast, the
corresponding section there is the place to slow down.

\textbf{The linear algebra book} supplies the mathematical language for the second half of
this volume. The instant we write down the Laplacian $\Lap = \mathbf{D} - \Adjm$
in the spectral chapters, we are doing spectral theory: eigenvalues,
eigenvectors, the Rayleigh quotient, positive semidefiniteness, the quadratic form
$x^\top \Lap x = \sum_{(u,v)\in E}(x_u - x_v)^2$. Every one of those is a concept
the linear algebra book builds from scratch. The whole spectral arc -- from the
Laplacian through Cheeger, expanders, random walks, and Laplacian solvers -- is,
in a real sense, an extended application of that book to graphs: a graph becomes a
matrix, a cut becomes a vector, a bottleneck becomes the second-smallest
eigenvalue, and ``solve this network'' becomes ``solve $\Lap x = b$.''
\end{connectionbox}

You can read this book on its own -- it is self-contained, and the appendix
gathers the data-structure and linear-algebra facts we lean on. But the Connection
boxes are written assuming the companions exist, and a reader who follows them gets
a unified picture instead of three disconnected courses.

\section*{The shape of the book: from BFS to Laplacian solvers}

Seeing the whole arc on one page before you start helps frame the road ahead,
because the later destinations explain the earlier choices. The book moves through five broad
territories.

\begin{focusbox}
\textbf{Territory I --- Foundations and traversal.} Representations, BFS, DFS,
topological order, strongly connected components, minimum spanning trees, and
shortest paths from Dijkstra up through Johnson's all-pairs algorithm. We re-pour
this foundation at full rigor: every greedy choice gets an exchange argument,
every relaxation gets an invariant.

\textbf{Territory II --- Cuts, flows, and matchings.} The duality-rich heart of
combinatorial optimization: bridges and biconnectivity, global minimum cut,
maximum flow and its min-cut dual, push-relabel, min-cost flow, bipartite and
weighted bipartite matching, Edmonds' blossom algorithm for general matching, and
the art of modeling a problem as a flow.

\textbf{Territory III --- Structure on trees and sparse graphs.} Heavy-light and
centroid decomposition, Euler tours, link-cut trees, LCA and range structures,
treewidth and tree decompositions, fixed-parameter tractability, planarity,
separators, coloring, independent sets and cliques, and perfect graphs. Here
structure is the resource: exploit it and intractable problems become easy.

\textbf{Territory IV --- The spectral world.} The Laplacian, spectral graph
theory, spectral clustering, expanders, the deep equivalence between random walks
and electrical networks, and PageRank as a Markov chain. This is where graphs
become linear algebra.

\textbf{Territory V --- Hardness, scale, and the frontier.} NP-hard graph
problems, approximation, dynamic graphs, streaming and massive graphs, and finally
near-linear-time Laplacian solvers and electrical flows -- the summit, where
Territory IV's spectral view and Territory II's flow view fuse into a single
result that ranks among the field's deepest.
\end{focusbox}

The thread that ties the summit to the start is a single quantity that mutates as
we climb. Connectivity, in the BFS chapter, is a yes/no question answered by BFS:
can I reach $t$ from $s$? In the maximum-flow chapter it has become a number:
\emph{how much} can flow from $s$ to $t$, equal by duality to the lightest cut
separating them. In the spectral chapters it has become spectral: how
well-connected is the graph \emph{as a whole}, measured by the second eigenvalue
$\lambda_2$ of the Laplacian, with Cheeger's inequality pinning that eigenvalue to
the sparsest cut. And in the random-walks chapter it is physical: connectivity is
the conductance of a resistor network, computable by solving a linear system. Same
question, four languages, each strictly more expressive than the last. Watching that
single idea grow up is the real story of the book, summarized across its four stages below:

\begin{center}
\begin{tikzpicture}[
    >=Stealth,
    stage/.style={draw=sourceblue!75!black, fill=sourceblue!6, rounded corners=4pt, thick,
      minimum width=3.15cm, minimum height=1.85cm, align=center, font=\footnotesize},
    arr/.style={->, very thick, draw=edgegray}
]
  \node[stage] (s1) at (0, 0) {
    \textbf{1. Traversal}\\[2pt]
    Reachable $s \leadsto t$?\\[2pt]
    {\scriptsize\color{sourceblue!90!black} BFS / DFS trees}
  };
  \node[stage, draw=targetgreen!75!black, fill=targetgreen!6] (s2) at (3.85, 0) {
    \textbf{2. Combinatorial}\\[2pt]
    Max $s$--$t$ flow $= $ min cut\\[2pt]
    {\scriptsize\color{targetgreen!85!black} Capacities \& cuts}
  };
  \node[stage, draw=quotientorange!80!black, fill=quotientorange!7] (s3) at (7.7, 0) {
    \textbf{3. Spectral}\\[2pt]
    Global bottleneck $\lambda_2(\Lap)$\\[2pt]
    {\scriptsize\color{quotientorange!85!black} Cheeger \& expanders}
  };
  \node[stage, draw=structurepurple!80!black, fill=structurepurple!7] (s4) at (11.55, 0) {
    \textbf{4. Electrical}\\[2pt]
    Effective resistance\\[2pt]
    {\scriptsize\color{structurepurple!85!black} Solve $\Lap x = b$ in $\widetilde{\Oh}(m)$}
  };

  \draw[arr] (s1.east) -- (s2.west);
  \draw[arr] (s2.east) -- (s3.west);
  \draw[arr] (s3.east) -- (s4.west);
\end{tikzpicture}
\end{center}

\section*{A worked trace, so you know what they look like}

The single habit this book most wants you to build is \emph{tracing an algorithm
by hand on a small graph}. Reading pseudocode tells you what an algorithm says;
tracing it tells you what it \emph{does}. So let us do one now, in the preface,
before any chapter -- both to set the standard and to dust off your graph
intuition. We will trace plain breadth-first search, the most elementary algorithm
in the book, because the point here is the \emph{method}, not the difficulty.

Consider this little undirected graph on six vertices. We will run BFS from $s$
and ask which vertices it reaches, in what order, and at what distance.

\begin{center}
\begin{tikzpicture}[
    every node/.style={circle,draw=nodeblue,fill=blue!6,thick,minimum size=8mm,
      font=\small},
    edge/.style={draw=edgegray,thick}]
  \node (s) at (0,0)   {$s$};
  \node (a) at (2,1)   {$a$};
  \node (b) at (2,-1)  {$b$};
  \node (c) at (4,1)   {$c$};
  \node (d) at (4,-1)  {$d$};
  \node (e) at (6,0)   {$e$};
  \draw[edge] (s) -- (a);
  \draw[edge] (s) -- (b);
  \draw[edge] (a) -- (c);
  \draw[edge] (b) -- (d);
  \draw[edge] (a) -- (b);
  \draw[edge] (c) -- (e);
  \draw[edge] (d) -- (e);
\end{tikzpicture}
\end{center}

\begin{intuitionbox}
BFS explores in rings. It visits everything at distance $0$ (just $s$), then
everything at distance $1$ (the neighbors of $s$), then distance $2$, and so on,
never touching a vertex at distance $k+1$ before all of distance $k$ is done. A
FIFO queue enforces exactly this order for free: a vertex enqueued earlier comes
out earlier, and since distance-$k$ vertices are all enqueued before any
distance-$(k+1)$ vertex, the rings stay separated. That outward ripple geometry ensures
that every vertex is reached along a path of minimum edge count.
\end{intuitionbox}

Here is the algorithm we are tracing, stated precisely.

\begin{algorithm}[H]
\caption{Breadth-First Search from a source}
\label{ch00:alg-bfs}
\SetKwInOut{Input}{Input}\SetKwInOut{Output}{Output}
\Input{Graph $\Gph=(V,E)$ as adjacency lists; source $s \in V$.}
\Output{$\dist[v]$ = number of edges on a shortest $s$--$v$ path, or $\infty$.}
\ForEach{$v \in V$}{$\dist[v] \gets \infty$}
$\dist[s] \gets 0$\;
$Q \gets$ empty FIFO queue; enqueue $s$\;
\While{$Q$ is not empty}{
  $u \gets$ dequeue($Q$)\;
  \ForEach{neighbor $w$ of $u$}{
    \If{$\dist[w] = \infty$}{
      $\dist[w] \gets \dist[u] + 1$\;
      enqueue $w$ into $Q$\;
    }
  }
}
\Return $\dist$\;
\end{algorithm}

Now the trace. Assume each adjacency list is sorted alphabetically, so neighbors
are scanned in that order -- pinning down the order makes the trace reproducible,
which is exactly what you want when debugging. We show the queue (front on the
left) and the distances discovered after each dequeue.

\begin{center}
\renewcommand{\arraystretch}{1.25}
\begin{tabular}{l l l}
\toprule
\textbf{Step} & \textbf{Queue (front $\to$ back)} & \textbf{Distances discovered this step} \\
\midrule
init        & $[\,s\,]$    & $\dist[s]=0$ \\
dequeue $s$ & $[\,a,b\,]$  & $\dist[a]=1,\ \dist[b]=1$ \\
dequeue $a$ & $[\,b,c\,]$  & $\dist[c]=2$ \ (neighbor $b$ already seen) \\
dequeue $b$ & $[\,c,d\,]$  & $\dist[d]=2$ \ (neighbor $a$ already seen) \\
dequeue $c$ & $[\,d,e\,]$  & $\dist[e]=3$ \\
dequeue $d$ & $[\,e\,]$    & (neighbor $e$ already seen) \\
dequeue $e$ & $[\ ]$       & (neighbors $c,d$ already seen) \\
\bottomrule
\end{tabular}
\end{center}

Redrawing the graph with vertices grouped into columns by distance turns the trace table into a layered BFS tree, separating the solid tree edges that first discover each vertex from the dashed cross edges between already-seen vertices:

\begin{center}
\begin{tikzpicture}[
    >=Stealth,
    layerbox/.style={draw=structurepurple!35, fill=structurepurple!4, rounded corners=3pt, dashed, inner sep=6pt},
    vnode/.style={circle, draw=sourceblue, fill=sourceblue!10, thick, minimum size=7.5mm, font=\small\bfseries},
    treeedge/.style={draw=sourceblue!85!black, very thick},
    crossedge/.style={draw=quotientorange!85!black, thick, dashed},
    labelstyle/.style={font=\footnotesize\scshape, text=gray!80!black}
]
  % Vertices partitioned into distance layers
  \node[vnode] (vs) at (0, 0) {$s$};
  \node[vnode] (va) at (2.4, 0.9) {$a$};
  \node[vnode] (vb) at (2.4, -0.9) {$b$};
  \node[vnode] (vc) at (4.8, 0.9) {$c$};
  \node[vnode] (vd) at (4.8, -0.9) {$d$};
  \node[vnode] (ve) at (7.2, 0) {$e$};

  % Layer bounding containers
  \begin{scope}[on background layer]
    \node[layerbox, fit=(vs), minimum width=1.5cm, minimum height=2.8cm] (L0) {};
    \node[layerbox, fit=(va) (vb), minimum width=1.5cm, minimum height=2.8cm] (L1) {};
    \node[layerbox, fit=(vc) (vd), minimum width=1.5cm, minimum height=2.8cm] (L2) {};
    \node[layerbox, fit=(ve), minimum width=1.5cm, minimum height=2.8cm] (L3) {};
  \end{scope}

  % Layer headers
  \node[labelstyle, above=0.25cm of L0.north] {$L_0$ ($\dist=0$)};
  \node[labelstyle, above=0.25cm of L1.north] {$L_1$ ($\dist=1$)};
  \node[labelstyle, above=0.25cm of L2.north] {$L_2$ ($\dist=2$)};
  \node[labelstyle, above=0.25cm of L3.north] {$L_3$ ($\dist=3$)};

  % Tree edges
  \draw[treeedge] (vs) -- (va);
  \draw[treeedge] (vs) -- (vb);
  \draw[treeedge] (va) -- (vc);
  \draw[treeedge] (vb) -- (vd);
  \draw[treeedge] (vc) -- (ve);

  % Non-tree / cross edges
  \draw[crossedge] (va) -- (vb) node[midway, right=3pt, font=\scriptsize, text=quotientorange!90!black] {cross};
  \draw[crossedge] (vd) -- (ve) node[midway, below right=2pt, font=\scriptsize, text=quotientorange!90!black] {cross};

  % Legend
  \node[anchor=center, font=\scriptsize, text=gray!80!black] at (3.6, -2.0) {
    \textbf{Classification:}\quad
    \tikz[baseline=-0.5ex]{\draw[treeedge] (0,0) -- (0.5,0);} BFS Tree Edge \qquad
    \tikz[baseline=-0.5ex]{\draw[crossedge] (0,0) -- (0.5,0);} Cross Edge
  };
\end{tikzpicture}
\end{center}

The final distances are $\dist[s]=0$, $\dist[a]=\dist[b]=1$,
$\dist[c]=\dist[d]=2$, $\dist[e]=3$. Read the rings straight off the queue: $s$
alone, then $\{a,b\}$, then $\{c,d\}$, then $\{e\}$. Notice the two edges that
discovered nothing -- the edge $a$--$b$ when we processed $a$, and the edge
$d$--$e$ when we processed $d$. Those are precisely the edges joining vertices in
the same ring or an already-finished ring; in the directed-graph chapters that
distinction between ``tree'' and ``non-tree'' edges becomes the engine behind
cycle detection and strongly connected components. A trace exposes the underlying
anatomy directly: it lays bare the structural boundary between traversed tree
branches and cross-connections.

\begin{whybox}
\emph{Why the distances are correct.} The intuition is that BFS explores outward in
concentric rings. Every vertex at distance $1$ is found directly from the source $s$;
every vertex at distance $2$ is found from some vertex at distance $1$, and so on.
Because vertices are processed in first-in, first-out order, all vertices at distance
$k$ are explored before any vertex at distance $k+1$. When a vertex $w$ is first
encountered from an active vertex $u$, the path from $s$ through $u$ to $w$ has
length $\dist[u] + 1$, which is as short as possible. We will prove this ring-by-ring
property carefully in Chapter~2 with a formal induction; here, the key takeaway is
how the FIFO queue structure directly drives the discovery order.
\end{whybox}

\begin{complexitybox}
Each vertex is enqueued and dequeued at most once -- it is enqueued only when its
distance flips from $\infty$, which happens once. So the \textbf{while} loop runs
$\Oh(n)$ times. Across all iterations the inner \textbf{foreach} scans each
adjacency list exactly once, totaling $\sum_v \deg(v) = 2m$ edge inspections in an
undirected graph. With adjacency lists and an array-backed queue, every individual
operation is $\Oh(1)$. Total: $\Th(n+m)$ time and $\Th(n)$ extra space. This is
optimal -- you cannot decide reachability without, in the worst case, looking at
every vertex and edge once.
\end{complexitybox}

\begin{trapbox}
The most common BFS bug, and a template for the bugs this book hunts everywhere:
marking a vertex as visited \emph{when you dequeue it} instead of \emph{when you
enqueue it}. In our trace, $b$ is a neighbor of both $s$ and $a$. With
enqueue-time marking, $b$ is discovered once (by $s$) and added to the queue once.
With dequeue-time marking, both $s$ and $a$ would enqueue $b$, so $b$ enters the
queue twice; you may still get correct distances, but the queue can swell to
$\Th(m)$ and on a dense graph your ``linear'' algorithm quietly degrades. The
guard \texttt{if}~$\dist[w]=\infty$ at enqueue time is doing real work. Lesson:
the invariant serves as the exact contract the implementation must satisfy at every step.
\end{trapbox}

That is the rhythm of the entire book, here on the gentlest possible example:
picture, pseudocode, trace, why-it-works, cost, trap. When the example is Edmonds'
blossom contraction or a Cheeger cut instead of BFS, the rhythm is identical and
the boxes mean the same things. Learn the rhythm once and the hard chapters become
\emph{long}, never \emph{foreign}.

\section*{How to read this book}

Read it with a pen. Every \emph{Trace it yourself} box is an invitation to run the
algorithm before reading our version, and the gap between your trace and ours is
the most precise feedback you will get about your own understanding. Draw the small
graphs; the act of placing the vertices is half the insight.

\begin{marginnotebox}
Two reading paths. \textbf{Linear} -- chapters are ordered so each rests only on
earlier ones, and reading straight through is the intended default.
\textbf{Targeted} -- if you came for a specific tool (say, max-flow for a project,
or Laplacian solvers for a paper), jump to its chapter and chase the backward
Connection boxes to fill gaps; they form a dependency trail you can follow
upstream. The final appendix is your safety net for any data-structure or
linear-algebra prerequisite invoked without ceremony in the main text.
\end{marginnotebox}

Do not skip the proofs, but do not let them block you either. On a first pass it
is legitimate to read a theorem, absorb its \emph{Why it works} box, run the
trace, and defer the line-by-line proof to a second pass. Understanding what is
true and why it should be true is enough to use an algorithm; understanding the
proof in full is what lets you modify it for the variant your problem actually
needs -- and sooner or later your problem will need a variant.

\begin{historybox}
A last word on why the order is what it is. The field did not develop top-down.
BFS and DFS were folklore made rigorous; Dijkstra (1956) and Kruskal/Prim gave
greedy its proudest victories; Ford and Fulkerson (1956) found max-flow/min-cut
and with it the duality that organizes half this book; Edmonds (1965) drew the
line between polynomial and not and invented blossoms to stay on the right side of
it; Tarjan (1970s) showed a single depth-first search could compute almost any
structural property of a graph; Fiedler (1973) connected a graph's connectivity to
its second eigenvalue; Karger (1990s) let randomness find minimum cuts; and
Spielman and Teng (2004) closed the loop by solving Laplacian systems in
near-linear time, making the spectral view computationally as cheap as the
combinatorial one. We follow roughly that historical arc because the problems came
before the solutions, and meeting the problem first is the surest way to
understand why anyone bothered to solve it.
\end{historybox}

\section*{Exercises}

These are warm-ups, not gatekeepers. They use nothing past this preface, and their
only purpose is to get a pen in your hand before Chapter~1. Solutions are not
printed; the point is the doing, and a trace you can defend to yourself is its own
answer key.

\begin{exercise}[Trace it cold]
Redraw the undirected six-vertex graph from the worked trace, then add one edge:
$c$--$d$. Re-run BFS from $s$ by hand, writing the queue and the discovered
distances after each dequeue exactly as we did. Which distances change, and which
do not? Identify, for the new edge, whether it is discovered as a tree edge or
found pointing at an already-seen vertex.
\end{exercise}

\begin{exercise}[Changing the starting point]
On the undirected graph from the worked trace, run the search starting from vertex
$c$ instead of $s$. In what order are the vertices discovered, and what is the
maximum distance from $c$ to any reachable vertex?
\end{exercise}

\begin{exercise}[Break the algorithm on purpose]
Implement the dequeue-time-marking bug described in the \emph{Common trap} box:
move the ``mark as visited'' step from enqueue time to dequeue time. On the
undirected graph, exhibit a vertex that enters the queue twice. Then construct, or
describe, a family of graphs on $n$ vertices where this bug pushes the queue to
size $\Th(n^2)$ even though the correct algorithm keeps it at $\Oh(n)$.
\end{exercise}

\begin{exercise}[Draw your own graph]
Draw an undirected connected graph with 6 vertices where every vertex has degree at
least 2, but the shortest-path distance between some pair of vertices is at least 3.
Trace a search from one endpoint of that diameter to the other.
\end{exercise}

\begin{exercise}[A counting sanity check]
The \emph{Cost summary} box asserts $\sum_v \deg(v) = 2m$ for an undirected graph.
Count the degrees of all six vertices in the worked undirected graph, sum them, and
confirm the total is twice the number of edges you drew. Explain in your own words
why every undirected edge contributes exactly 2 to the sum of degrees.
\end{exercise}

The summit -- a near-linear-time Laplacian solver computing electrical flows -- is
demanding, and we will not pretend otherwise. But by the time you reach it
you will have built every piece: the graph as a matrix from the Laplacian chapter,
the spectrum and Cheeger's inequality from the spectral and expander chapters, the
random-walk and electrical correspondence from the random-walks chapter, the flow
duality from the maximum-flow chapter, and the sparsification ideas that let a
huge linear system be solved by repeatedly solving small ones. Nothing at the summit is
built from scratch; it assembles the exact tools you developed on the way up.

Turn the page. Before any algorithm can traverse a graph, we must answer two
basic questions: what is a graph mathematically, and how do we store its vertices
and edges in computer memory so that looking up neighbors is fast?
In Chapter~1, we lay this groundwork by formalizing directed and undirected graphs
and evaluating adjacency lists, adjacency matrices, and edge lists under our asymptotic
cost model.

\chapterline
